\sezione{11}{Esempio maschi/femmine con Bayes parametrico}{54–57}

Lo \textbf{stesso} problema della sezione 4, risolto con le gaussiane invece del conteggio nel cerchio. Obiettivo: classificare $\x=[57,\,168]\T$ (peso in kg, altezza in cm).

\subsection{I dati}
\begin{center}\small
\begin{tabular}{cc}
\begin{tabular}{ccc}
\rowcolor{blu}\thd{Peso} & \thd{Altezza} & \thd{Classe}\\
\rowcolor{white}72&173&$w_1$\\ \rowcolor{blu!8}54&159&$w_1$\\ 65&172&$w_1$\\ \rowcolor{blu!8}58&170&$w_1$\\
62&165&$w_1$\\ \rowcolor{blu!8}72&176&$w_1$\\ 60&173&$w_1$\\ \rowcolor{blu!8}64&179&$w_1$\\
\end{tabular}
&
\begin{tabular}{ccc}
\rowcolor{rosso}\thd{Peso} & \thd{Altezza} & \thd{Classe}\\
\rowcolor{white}55&166&$w_2$\\ \rowcolor{rosso!8}46&158&$w_2$\\ 52&158&$w_2$\\ \rowcolor{rosso!8}47&160&$w_2$\\
55&167&$w_2$\\ \rowcolor{rosso!8}54&166&$w_2$\\ 55&164&$w_2$\\ \rowcolor{rosso!8}49&157&$w_2$\\
51&165&$w_2$\\ \rowcolor{rosso!8}51&162&$w_2$\\
\end{tabular}
\end{tabular}
\end{center}

\subsection{Passo 1: stima dei parametri (MLE, divisione per $n$)}
\begin{esempio}
\[
\mmu_1=[63.4,\ 170.9]\T\qquad \SSigma_1=\begin{bmatrix}35.2&23.3\\23.3&34.9\end{bmatrix}
\qquad\qquad
\mmu_2=[51.5,\ 162.3]\T\qquad \SSigma_2=\begin{bmatrix}10.1&8.9\\8.9&13.0\end{bmatrix}
\]
Entrambe le covarianze hanno $\sigma^{12}>0$ (peso e altezza crescono insieme). La nuvola delle femmine è molto più \textbf{concentrata}: $|\SSigma_2|\approx52.4$ contro $|\SSigma_1|\approx685.5$.
\end{esempio}

\subsection{Passo 2: likelihood, evidence, posteriori}
\begin{esempio}[Prior come nella sezione 4: $P(w_1)=8/18$, $P(w_2)=10/18$]
\[
p(\x|w_1)=\frac{1}{(2\pi)^{d/2}|\SSigma_1|^{1/2}}\exp\!\left[-\tfrac12(\x-\mmu_1)\T\SSigma_1^{-1}(\x-\mmu_1)\right]\approx0.0033
\]
\[
p(\x|w_2)=\frac{1}{(2\pi)^{d/2}|\SSigma_2|^{1/2}}\exp\!\left[-\tfrac12(\x-\mmu_2)\T\SSigma_2^{-1}(\x-\mmu_2)\right]\approx0.0045
\]
\[
p(\x)=0.0033\cdot\tfrac{8}{18}+0.0045\cdot\tfrac{10}{18}\approx0.0040
\]
\begin{center}
\begin{tcolorbox}[enhanced,width=.8\linewidth,colback=white,colframe=rosso,boxrule=1.2pt,arc=4pt,halign=flush center]
$P(w_1|\x)\approx0.36\qquad\qquad P(w_2|\x)\approx\mathbf{0.64}\quad\Rightarrow\quad\text{\textbf{femmina}}$
\end{tcolorbox}
\end{center}
(Con più cifre: $p(\x|w_2)=0.00455$; la slide riporta $0.0045$ troncando.)
\end{esempio}

\fig[.98\linewidth]{mf_param.pdf}
\didascalia{Sinistra: le due gaussiane stimate. Destra: la posteriori su tutto il piano e il confine di decisione, una curva perché $\SSigma_1\neq\SSigma_2$ (sezione 10). Corrisponde alla slide 57.}

\subsection{Un'osservazione sorprendente}
\begin{intuizione}[$\x$ è più vicino ai maschi, eppure vince «femmina»]
Calcolando la distanza di Mahalanobis da ciascuna classe:
\[
r_1^2\approx1.25\ (r_1\approx1.12)\qquad\qquad r_2^2\approx3.15\ (r_2\approx1.77)
\]
Rispetto alla propria nuvola, $\x$ è un maschio \emph{più tipico} di quanto sia una femmina tipica (anche in distanza euclidea: 7.0 contro 7.9). Perché allora vince $w_2$?
\begin{itemize}
  \item La gaussiana delle femmine è molto più \textbf{stretta}, quindi molto più \textbf{alta}: $|\SSigma_2|^{1/2}\approx7.2$ contro $|\SSigma_1|^{1/2}\approx26.2$. Il picco femminile è circa \textbf{3.6 volte} più alto.
  \item Il fattore esponenziale favorisce i maschi ($e^{-0.62}\approx0.54$ contro $e^{-1.58}\approx0.21$, cioè 2.6 volte), ma non abbastanza da compensare.
  \item Il prior (10/18 contro 8/18) dà un'ulteriore spinta a $w_2$. Anche con prior uguali vincerebbe comunque $w_2$, con $P(w_2|\x)\approx0.58$.
\end{itemize}
\textbf{Morale:} la likelihood gaussiana non è solo «quanto sei vicino al centro», ma anche «quanto è concentrata la classe». È il termine $-\frac12\ln|\SSigma_i|$ della funzione discriminante (sezione 10).
\end{intuizione}

\subsection{Confronto con la stima a conteggio}
\begin{center}
\rowcolors{2}{chiaro}{white}
\begin{tabularx}{\linewidth}{l c Y}
\rowcolor{navy}\thd{Metodo} & \thd{$P(w_2|\x)$} & \thd{Commento}\\
conteggio nel cerchio (sez. 4) & $2/3\approx0.67$ & dipende dal raggio, usa solo 3 punti\\
gaussiane (sez. 11) & $\approx0.64$ & usa tutti i 18 punti tramite $\mmu_i,\SSigma_i$; vale ovunque nel piano\\
\end{tabularx}
\end{center}
Le due stime concordano, ma la versione parametrica produce un \textbf{confine liscio} su tutto il piano e non dipende da una scelta arbitraria del raggio.
\begin{attenzione}
Con soli 8 e 10 esempi stiamo stimando 5 parametri per classe: le stime sono incerte (per esempio, con la divisione per $n-1$ le covarianze crescerebbero di circa l’11–14\%). Il risultato va letto come «probabilmente femmina», non come certezza.
\end{attenzione}

\begin{appunti}{— Sezione 11}
\pag{54} Stesso dataset della sezione 4 (8 M, 10 F); $\x=[57,168]\T$.
\pag{55} $\mmu_1=[63.4,170.9]$, $\SSigma_1=\begin{psmallmatrix}35.2&23.3\\23.3&34.9\end{psmallmatrix}$; $\mmu_2=[51.5,162.3]$, $\SSigma_2=\begin{psmallmatrix}10.1&8.9\\8.9&13.0\end{psmallmatrix}$ (MLE con $1/n$). Femmine più concentrate: $|\SSigma_2|\approx52$ vs $686$.
\pag{56} $p(\x|w_1)\approx0.0033$, $p(\x|w_2)\approx0.0045$, $p(\x)\approx0.0040$; $P(w_1|\x)\approx0.36$, $P(w_2|\x)\approx0.64$ → femmina.
\pag{57} Confine curvo ($\SSigma$ diverse). $\x$ più vicino ai maschi in Mahalanobis ($r^2$: 1.25 vs 3.15) ma vince F perché la sua gaussiana è $\approx3.6$ volte più alta + prior. Prior uguali: 0.58. Concorda con il conteggio (0.67).
\end{appunti}
